\chapter{Scenario layer parameters}
\label{app:scenario}

The scenario layer attaches a synthetic gym to the 4,000 real members, with a documented ground truth (\cref{sec:scenario}). Everything below is regenerated by the data build and is deterministic in the seed (42 for every reported result). Experiments that use it test whether the system recovers known structure and decides well against a known truth; they are not claims about how real members respond.

\section{Tables and parameters}

{\small\sffamily
\begin{xltabular}{\linewidth}{@{}l l L l@{}}
  \caption{The data tables written by the build: the hand-over from data to the rest of the system.}\label{tab:datatables}\\
  \toprule
  \hd{Table} & \hd{One row per} & \hd{Main columns} & \hd{Source} \\
  \midrule
  \endfirsthead
  \toprule
  \hd{Table} & \hd{One row per} & \hd{Main columns} & \hd{Source} \\
  \midrule
  \endhead
  \code{members} & member & identifier, the 14 Kaggle columns, club, monthly fee, split & Kaggle + synthetic \\
  \code{instructors}, \code{classes} & instructor, class & club, name; a class's instructor & synthetic \\
  \code{attendance} & member and class & weekly visits (always above 0) & synthetic \\
  \code{offers} & offer & cost, staff minutes, stock, channels & synthetic \\
  \code{interventions} & past contact & week, offer, channel, treated, stayed & synthetic \\
  \code{events} & gym event & type, entity, week & synthetic \\
  \code{true_uplift} & member and offer & uplift & ground truth \\
  \code{true_churn} & member & hazard before and after events, exposure shares & ground truth \\
  \code{features} & member & the 17 model features & derived \\
  \code{segments} & segment & name, description, size, churn rate, centre & model output \\
  \code{member_scores} & member & segment, churn probability, band, fee, club, top drivers & model output \\
  \code{uplift_tlearner} & member and offer & estimated uplift & model output \\
  \bottomrule
\end{xltabular}}

\begin{table}[tb]
  \caption{Generation parameters.}
  \label{tab:scenparams}
  \begin{tabularx}{\linewidth}{@{}l L@{}}
    \toprule
    \hd{Parameter} & \hd{Value} \\
    \midrule
    Scoring date & 12 October 2026 \\
    Clubs & C1, C2, C3; each member assigned uniformly at random \\
    Instructors and classes & 3 instructors per club, 2 classes each: 9 instructors, 18 classes \\
    Classes per class-going member & 1 to 3, uniformly, at the member's own club; 1,611 members attend \\
    Monthly fee & \$72, \$54 or \$46 for 1-, 6- or 12-month contracts \\
    History & 12 weeks; 60\% of members contacted once (2,400); 60\% of those treated in expectation (1,416) \\
    Event effects & 1.4 (instructor left) and 0.9 (class cancelled) in log-odds per full share of visits \\
    Seed & 42 \\
    \bottomrule
  \end{tabularx}
\end{table}

The instructors are Sam Rivera, Priya Nair and Lee Wong (C1); Maya Chen, Omar Haddad and Ana Costa (C2); and Tom Becker, Nora Ali and Ken Sato (C3). The classes are Spin, HIIT, Yoga, Pilates, Boxfit, Circuit, Zumba, Core, Stretch, Bootcamp, Barre, Step, Kettlebell, Aqua, Mobility, TRX, Dance and Power, two per instructor in that order. A class-going member's weekly visits are split across their classes with random weights so that the total equals their real Kaggle frequency.

\section{Ground truth}

\paragraph{Churn without contact.} Each member's hazard for the next month uses features only, never the real label:
\begin{equation}
  z = -1.4 + 1.3\,[\text{lifetime} \le 2] + 1.1\,[\text{trend} < 0.7] - 0.9\,[\text{contract} = 12] - 0.4\,\text{group} - 0.5\,(\text{visits} - 1.9),
\end{equation}
with base hazard $\sigma(z)$. It ranges from 0.01 to 0.87 and correlates 0.715 with the real churn label (\cref{fig:hazard}), so the scenario's risk is consistent with the real data without being fitted to it.

\paragraph{Offer effects.} With \emph{falling} = trend below 0.85, \emph{new} = lifetime at most 3 months and \emph{price-sensitive} = a monthly contract with extra spend below 120, the raw effects on next-month retention are: check-in $0.02 + 0.04\,\text{falling}$; free PT $0.02 + 0.14\,\text{new} + 0.04\,\text{falling}$; class pass $0.01 + 0.10\,\text{group} + 0.04\,\text{referred}$; discount $0.01 + 0.12\,\text{price-sensitive} - 0.06\,[\text{contract}=12]$. Each is multiplied by $(0.6 + \text{hazard})$, so persuasion matters more for members at risk, and the realised uplift saturates near certain retention: $\text{clip}(1 - h + e, 0, 1) - (1 - h)$. The discount for annual members is the only negative effect, the \enquote{sleeping dogs} case. Over all 16,000 pairs the mean true uplift is +5.8 points; \cref{fig:scenario} shows it by segment.

\paragraph{Gym events.} In the week of 5 October 2026 the instructor with the most attendance (Maya Chen, Zumba and Core, C2) leaves and the busiest class at another club (HIIT, C1) is cancelled. A member's event shift is $1.4\,s_\text{instructor} + 0.9\,s_\text{class}$ in log-odds, where $s$ are the shares of their class visits affected: a member who trained only with the departed instructor has their odds of leaving multiplied by $e^{1.4} \approx 4.1$. An event counts from its week; earlier contacts keep their pre-event outcomes. For the 186 members with at least half their visits with the departed instructor the hazard rises from 0.262 to 0.472 (×1.80); for the 170 attending the cancelled class from 0.253 to 0.334 (×1.32). Across all members it rises from 0.320 to 0.335.

\begin{figure}[tb]
  \centering
  \pgfplotstableread{\results{a5_event_sensitivity.csv}}\sens
  \includegraphics[width=0.46\linewidth]{generated/base_hazard.pdf}\hfill
  \begin{tikzpicture}
    \begin{axis}[anchor, scale only axis, width=0.36\linewidth, height=3.8cm, xlabel={scale of both event effects},
        ylabel={lift for lost-instructor members}, title={Event sensitivity}, xtick={0,0.5,1,1.5,2},
        legend style={at={(0.03,0.97)}, anchor=north west},
        xticklabel style={/pgf/number format/fixed, /pgf/number format/precision=1}]
      \addplot[cBlue, mark=*, mark size=1.5pt] table[x=scale, y=within_rise]{\sens};
      \addplot[cOrange, mark=square*, mark size=1.4pt] table[x=scale, y=lift_vs_all]{\sens};
      \draw[muted, dashed] (axis cs:1,0.7) -- (axis cs:1,2.8);
      \legend{within the same members, against all members}
    \end{axis}
  \end{tikzpicture}
  \caption{Left: the synthetic no-contact hazard over the 4,000 members. Right: how the lost-instructor signal depends on the assumed event effects; at scale 0 the selection of class-goers alone puts flagged members below average (0.82). The default scale is dashed.}
  \label{fig:hazard}
  \source{\code{scripts/figures.py} (\code{hazard}); \code{results/a5_event_sensitivity.csv}.}
\end{figure}

\paragraph{Intervention history.} 2,400 members were each contacted once in a random week of the 12 before the scoring date; each was treated with probability 0.6 with a uniformly random offer through a random allowed channel, and the outcome was drawn from $1 - h + u$ for treated members (with the hazard in force that week). Treated members stay 4.9 points more often than controls (\cref{fig:scenario}).

\section{Limitations of the scenario}

\begin{itemize}
  \item \textbf{The effects are assumptions} encoding plausible playbook knowledge, not measurements; recovering them shows that a model \emph{can} recover heterogeneous effects, not that these are the real ones.
  \item \textbf{Uplift grows with risk by design}, so ranking by risk within an offer already tracks the true uplift for three of four offers (Spearman 0.83, 0.87 and 0.68; 0.07 for the class pass), whereas field data shows the riskiest customers are not necessarily the most responsive~\cite{ascarza2018futility}.
  \item \textbf{No contact fatigue and no channel effect}, so guardrail G01 and the channel choice cannot show a benefit against this truth (E6).
  \item \textbf{A small, noisy history}: 2,400 yes-or-no outcomes with effects of a few points, enough to recover rankings against the truth but not to confirm them from outcomes alone (E3).
  \item \textbf{Only generated events are scored}; events recorded later in the application change the graph and rules but not the ground truth, and only the event effects were varied for sensitivity, not the offer effects.
\end{itemize}
